% Corrected French for printed p.144.
% Editorial qualification: sheaf-module categories need not lack projectives.
% Corrections: peut être considérée; homomorphismes de bord (twice).
\SourcePage{144}

\LayerTarget{tohoku:construction:derived-bifunctor:p2}
\noindent foncteurs dérivés partiels droits de $T$ par rapport à la seconde
variable $B$,
\[
 R_2^iT(A,B)=H^i(T(A,C(B)))
\]
où $C(B)$ est le complexe défini par une résolution droite de $B$ par des
objets injectifs. Bien entendu, les $R_2^iT$ sont des bifoncteurs.

\LayerTarget{tohoku:construction:derived-bifunctor:cohomological-first-variable}
\textit{Supposons maintenant que pour tout objet injectif $B$ dans $C_2$, le
foncteur $A\longrightarrow T(A,B)$ sur $C_1$ soit exact}, nous allons montrer
que pour tout $B\in C_2$ la suite $(R_2^iT(A,B))$ peut être considérée comme un
foncteur cohomologique en $A$. Soit en effet $C(B)$ le complexe défini par une
résolution droite de $B$ par des objets injectifs. Pour toute suite exacte
$0\longrightarrow A'\longrightarrow A\longrightarrow A''\longrightarrow0$
dans $C_1$, la suite d'homomorphismes de complexes
\[
0\longrightarrow T(A'',C(B))\longrightarrow T(A,C(B))
 \longrightarrow T(A',C(B))\longrightarrow0
\]
est exacte (d'après l'hypothèse sur $T$, les termes de $C(B)$ étant injectifs),
elle définit donc une suite exacte de cohomologie, c'est-à-dire une suite
exacte
\LayerTarget{tohoku:sequence:partial-derived-long-exact}
\[
\cdots\longrightarrow R_2^iT(A'',B)\longrightarrow R_2^iT(A,B)
\longrightarrow R_2^iT(A',B)
\mathrel{\mathop{\longrightarrow}\limits^{\partial}}R_2^{i+1}T(A'',B)
\longrightarrow\cdots .
\]
On vérifie aussitôt que le morphisme $\partial$ dans cette suite ne dépend pas
du choix particulier du complexe $C(B)$, et qu'il permute aux homomorphismes
de suites exactes, ce qui montre bien que $(R_2^iT(A,B))$ pour $B$ fixé est un
foncteur cohomologique en $A$ ; on vérifie d'ailleurs aussitôt que pour un
morphisme $B\longrightarrow B'$ dans $C_2$, les morphismes correspondants
$R_2^iT(A,B)\longrightarrow R_2^iT(A,B')$ définissent un morphisme de
$\partial$-foncteurs en $A$.

\LayerTarget{tohoku:def:right-balanced}
Si on suppose que dans $C_1$, tout objet est isomorphe à un quotient d'un objet
projectif, et que pour tout objet projectif $A\in C_1$, le foncteur $T(A,B)$
est exact en $B$ (on dit alors, conformément à la terminologie de [6], que $T$
est ``right balanced''), alors les $R_2^iT(A,B)$ sont aussi les foncteurs
dérivés droits $R^iT(A,B)$ de $T$, et coïncident aussi avec les foncteurs
dérivés partiels $R_1^iT(A,B)$ ; dans ce cas, les homomorphismes de bord
construits ci-dessus ne sont autres que les homomorphismes de bord naturels de
$R^iT$ et $R_1^iT$ par rapport à leur première variable.

\LayerTarget{tohoku:example:ext-bifunctor}
Les considérations précédentes ont été développées surtout pour le cas d'une
catégorie abélienne $C$ dans laquelle tout objet est isomorphe à un sous-truc
d'un objet injectif, mais où on ne suppose pas que tout objet soit isomorphe à
un truc quotient d'un objet projectif. Prenons alors $C_1=C_2=C$, $C'=$
catégorie des groupes abéliens, $T(A,B)=\mathrm{Hom}(A,B)$ ; par définition
même d'un objet injectif, $\mathrm{Hom}(A,B)$ est exact en $A$ si $B$ est
injectif. Désignant par $\mathrm{Ext}^p(A,B)$ les foncteurs dérivés droits par
rapport à $B$, on voit que les $\mathrm{Ext}^p(A,B)$ forment un foncteur
cohomologique en $B$ et en $A$. On a
$\mathrm{Ext}^0(A,B)=\mathrm{Hom}(A,B)$ puisque $\mathrm{Hom}(A,B)$ est exact
à gauche, il est facile de voir que $\mathrm{Ext}^1(A,B)$ peut encore
s'interpréter comme le groupe des classes d'extensions de $A$ (quotient) par
$B$ (sous-truc) [6 ; Chap. XIV].

\LayerTarget{tohoku:example:sheaf-ext}
Un cas important où la catégorie abélienne $C$ contient suffisamment d'objets
injectifs et peut ne pas contenir assez d'objets projectifs est fourni par les
catégories de faisceaux de modules sur un faisceau d'anneaux donné sur un
espace topologique $X$. Au Chap. IV, nous étudierons de façon plus détaillée
les groupes $\mathrm{Ext}^p(A,B)$ dans ce cas.

\smallskip
\noindent\textbf{Correction éditoriale.} La formulation imprimée affirme sans
restriction l'absence d'assez de projectifs. Elle admet des exceptions : si
$X$ est réduit à un point, la catégorie est une catégorie de modules, qui a
assez de projectifs.

\LayerTarget{tohoku:section:2.4}
\noindent\textbf{2.4. Suites spectrales et foncteurs spectraux.} Pour la
théorie des

\TohokuContinuation{145}
